% Source: https://artofproblemsolving.com/community/c3h1225016_geometry
% Stable identity: aops-topic-1225016-q01
\emph{Step 1.} For any point $X$ of the triangle and any point $P$, the distance $PX$ is at most the largest of $PA$, $PB$, $PC$. Indeed $X$ lies on a segment $AY$ with $Y$ on side $BC$; $PY\le\max(PB,PC)$ because the distance from $P$ to points of a segment is largest at an endpoint (it is a convex function along the segment), and similarly $PX\le\max(PA,PY)$.

\emph{Step 2.} Applying Step 1 with $P=M$: $MN\le\max(MA,MB,MC)$. Applying it again to each of these, with the roles reversed: $MA\le\max(AB,AC,AA)$, etc. So $MN\le$ the largest distance between two vertices, which is $BC$.

\begin{asy}
usepackage("amsmath");
size(3.5cm);
defaultpen(fontsize(10pt));
pair A=(1,3),B=(0,0),C=(5,0),M=(1.2,1.5),N=(3.3,0.8);
draw(A--B--C--cycle,linewidth(1)); draw(M--N,blue); dot(M);dot(N);
label("$A$",A,N);label("$B$",B,SW);label("$C$",C,SE);label("$M$",M,W);label("$N$",N,E);
label("\mbox{}",(min(currentpicture,user=true)+max(currentpicture,user=true))/2,black); // keeps the text after the figure black when LaTeX typesets the labels inline
\end{asy}

Equality requires equality in each step, which forces $M$ and $N$ to be vertices at distance $BC$, i.e. $\{M,N\}=\{B,C\}$ (if $BC$ is strictly the longest side). In particular, for $M,N$ strictly inside, $MN<BC$. $\blacksquare$
